\section{Architecture}

The application is structured into three modules.
Two additional modules serve to test the agent and document the application.

\subsection{Frontend}

Spartan Write's frontend is built with React.js \cite{react} on Tauri, enabling cross-platform development of native desktop applications for macOS, Windows, and Linux from a single TypeScript/React codebase \cite{tauri}.

\subsubsection{Development tooling and libraries}

The frontend leverages modern development tooling to ensure productivity and maintainability while maximizing accessiblity.

\begin{itemize}
    \item Tauri provides seamless integration with native system APIs while minimizing binary size, making it ideal for resource-conscious deployment across diverse user environments \cite{tauri}.
    \item Vite powers fast build times and hot module reloading during development, as well as a lightweight packaging pipeline \cite{vite}.
    \item TypeScript provides type safety across the codebase \cite{typescript}.
    \item Styling is handled by Tailwind CSS for responsive, utility-first design \cite{tailwind}.
    \item User interface components are derived from shadcn, a plug-and-play component library designed with accessibility, extensibility, and consistent visual design in mind \cite{shadcn}.
    \item The source code editor implements the Monaco Editor, which provides syntax-aware code editing with support for \LaTeX\ and other file formats \cite{monaco}.
\end{itemize}

\subsubsection{Agentic AI chat integration}

A key feature of the frontend is its integration with CopilotKit, which enables agentic chat directly within the editor \cite{copilotkit}.
The agentic chat system supports tool use for file operations, \LaTeX\ compilation, and image handling.
This tight integration allows users to request document modifications conversationally without leaving the editor.

\subsubsection{Dual-mode editor layout}

The editor employs a resizable split-panel layout that supports two complementary viewing modes.
In \emph{source view}, users see a syntax-highlighted code editor alongside a file browser for project navigation.
In \emph{preview view}, which is the default view, users see a live PDF rendering of their compiled document.
This dual-mode design provides immediate visual feedback and allows users to verify that AI-generated changes produce the desired output.

\subsubsection{Accessible interface}

Accessibility is a core design principle.
The interface includes configurable text sizing to accommodate users with vision impairments, screen reader support with live region announcements for screen reader users, a magnifier mode for enhanced visibility, and customizable PDF zoom levels ranging from 50\% to 300\%.

\subsubsection{Authentication}

The frontend integrates WorkOS AuthKit to provide secure, standards-based authentication \cite{workos}.
Users sign in via an external browser or through OAuth deeplinks, and the system manages JWT-based token refresh to maintain active sessions.
The authentication flow is designed to be seamless while adhering to security best practices.
Once authenticated, the frontend manages user credentials and provides visualization of usage data and costs.

\subsubsection{Project management}

The frontend provides a comprehensive project management interface.
Users access a dashboard to view and open recent projects.
There is an onboarding flow for new users, and a settings page displaying user profile information and usage analytics.
Native file dialogs (via Tauri plugins) allow users to select or create projects directly from the filesystem.

% TODO
\subsection{Sidecar}

\begin{itemize}
    \item Frontend is packaged along with a Python FastAPI server that will run alongside the application on the user's device.
    \item This server contains all of the logic to interact with the local file system.
    \item This allows headless usage of the application via HTTP.
\end{itemize}

\subsection{Agent}

\begin{itemize}
    \item Agent is a Python FastAPI server deployed to a trusted execution environment (TEE) on the cloud.
    \item Defines the behaviour of the AI agent using LangGraph.
    \item Uses a configurable LLM provision aggregator called OpenRouter, allowing for flexibility in the selection of the language model and the provider.
    \item After validating the authorization of the user, it invokes the agent, and handles subsequent interactions for frontend tool invocation and passing information.
\end{itemize}

\subsection{Benchmark}

\begin{itemize}
    \item The benchmark directly interacts with the agent module.
    \item Uses a curated dataset of document states along with test user prompts.
    \item Each item is scored on various subjective and objective metrics.
\end{itemize}

\subsection{User guide}

\begin{itemize}
    \item Documentation server.
\end{itemize}
